\section{Middle of a Linked List}
\label{sec:ll-middle}

\problemstatement{Given the head of a singly linked list, return its middle
node in a single pass. For even length, return the second of the two middle
nodes.}

\begin{patternbox}
\emph{``Find the middle in one pass''} is the \textbf{slow/fast two-pointer}
(tortoise--hare) technique: move one pointer one step and another two steps;
when the fast pointer reaches the end, the slow pointer is at the middle.
\end{patternbox}

\subsection*{Two speeds, one pass}

Start both pointers at the head. Each iteration advances \texttt{slow} by
one node and \texttt{fast} by two. When \texttt{fast} (or
\texttt{fast->next}) runs off the end, \texttt{slow} has travelled exactly
half as far -- it sits at the middle. This avoids first computing the length
and then re-walking, doing everything in one traversal.

\begin{ideabox}[Half Speed Lands on the Middle]
Because \texttt{fast} covers twice the distance of \texttt{slow}, by the
time \texttt{fast} finishes the list \texttt{slow} has covered half of it.
The loop condition (\texttt{fast \&\& fast->next}) chooses the second middle
for even lengths.
\end{ideabox}

\subsection*{Algorithm}

\begin{enumerate}
  \item \texttt{slow = fast = head}.
  \item While \texttt{fast} and \texttt{fast->next} are non-null: advance
        \texttt{slow} by one, \texttt{fast} by two.
  \item Return \texttt{slow}.
\end{enumerate}

\subsection*{Dry run}

\dryrun{List $1\to2\to3\to4\to5$.}

\begin{itemize}
  \item slow$=1$,fast$=1$. Step: slow$=2$,fast$=3$. Step: slow$=3$,
        fast$=5$.
  \item fast->next is null, stop. Middle $=3$.
\end{itemize}

\subsection*{C++ implementation}

\begin{lstlisting}
#include <bits/stdc++.h>
using namespace std;

struct ListNode {
    int val; ListNode* next;
    ListNode(int v) : val(v), next(nullptr) {}
};

ListNode* build(const vector<int>& a) {
    ListNode dummy(0), *t = &dummy;
    for (int x : a) { t->next = new ListNode(x); t = t->next; }
    return dummy.next;
}

ListNode* middleNode(ListNode* head) {
    ListNode* slow = head;
    ListNode* fast = head;
    while (fast && fast->next) {
        slow = slow->next;
        fast = fast->next->next;
    }
    return slow;
}

int main() {
    ListNode* head = build({1, 2, 3, 4, 5});
    cout << middleNode(head)->val << "\n";        // 3
    ListNode* even = build({1, 2, 3, 4});
    cout << middleNode(even)->val << "\n";         // 3 (second middle)
    return 0;
}
\end{lstlisting}

\begin{complexitybox}
\textbf{Time Complexity.} $O(n)$ -- one pass. \textbf{Space Complexity.}
$O(1)$.
\end{complexitybox}

\begin{takeawaybox}
The slow/fast two-pointer finds the middle in one pass and is the single
most reused linked-list trick: cycle detection, $n$-th-from-end, and
palindrome checks all start from it. Learn its even/odd termination
precisely.
\end{takeawaybox}
